% =========================================
% CHAPTER 3: RELATED WORK
% =========================================
\chapter{Related Work}
\label{ch:relatedwork}

This chapter surveys the academic literature and industry implementations that inform the comparative evaluation, then positions this thesis within the broader research landscape. The first section reviews key research contributions on rollup architectures. The second section examines prominent production systems on Ethereum mainnet. The third section analyzes gaps in the existing knowledge and describes how this thesis addresses each gap through the implemented framework.

\section{Academic Research on Layer~2 Protocols}

\subsection{Systematization of Knowledge}

Thibault et al. provide the most thorough systematization of Layer~2 blockchain protocols~\cite{sok-layer2}. Their taxonomy organizes the design space into four families. Payment channels, exemplified by the Lightning Network and Raiden, enable bidirectional off-chain transfers between fixed pairs of participants. Sidechains, including Plasma and various commit-chain designs, operate as semi-independent chains that periodically anchor their state to the parent chain. Rollups, encompassing both ZK and Optimistic variants, post transaction data on-chain while moving execution off-chain. Hybrid approaches such as Validium and Volition blend on-chain and off-chain data availability to offer users a configurable trade-off between cost and security.

This systematization provides a taxonomy and analyzes security properties at a theoretical level, but does not include empirical measurement of how these systems behave under failure conditions. The absence of controlled experimentation leaves open the question of whether theoretical security guarantees translate into practical resilience.

\subsection{ZK Rollup Research}

Bunz et al. introduced Bulletproofs, a proof system that generates efficient zero-knowledge range proofs without requiring a trusted setup ceremony~\cite{zk-efficient-proofs}. By eliminating the trusted setup, Bulletproofs removed a significant barrier to deploying zero-knowledge schemes in decentralized settings. Subsequent research extended this line of work toward SNARKs capable of verifying general computation, enabling ZK~Rollups that process arbitrary smart contract logic on Layer~2.

Gluchowski et al. published the architectural design of zkSync, a production-grade ZK~Rollup built on a PLONK-based proof system~\cite{zksync-paper}. Their contribution centers on proof generation optimizations and strategies for achieving compatibility with the Ethereum Virtual Machine. The Scroll project builds on this lineage with its Groth16-based verifier~\cite{scroll2023}. The \texttt{ZKVerifier} contract in this thesis draws on patterns from both zkSync and Scroll verification architectures.

Neither body of work provides a systematic evaluation of ZK~Rollup finality behavior under adverse conditions such as Layer~1 chain reorganizations or sequencer outages. The performance characteristics of ZK~Rollups in the presence of failure remain largely uncharted territory.

\subsection{Optimistic Rollup Research}

Kalodner et al. introduced Arbitrum, one of the earliest Optimistic~Rollups~\cite{arbitrum-original}. Their principal contribution is the interactive fault game, a multi-round dispute protocol that progressively narrows the scope of disagreement until the Layer~1 contract needs to re-execute only a single instruction to determine whether the challenged state transition was valid. This approach reduces on-chain cost of dispute resolution compared to re-executing the entire batch.

The \texttt{OptimisticVerifier} contract in this thesis incorporates patterns from Optimism Bedrock~\cite{optimism2023bedrock} and Arbitrum Nitro~\cite{nitro2023}, including a configurable challenge period through which commitments must pass before reaching finality. The implementation uses a period of ten blocks, a compressed value that preserves the architectural semantics of production systems while keeping experimental turnaround times manageable.

Kalodner et al. discuss sequencer liveness and censorship resistance in their design documents, but these properties have not been empirically evaluated in a controlled setting. The costs associated with failure recovery mechanisms remain largely unquantified in the existing literature.

\section{Industry Implementations}

Several production Layer~2 systems are deployed on Ethereum mainnet. This section reviews the most prominent ones as context for the design choices in the framework. Each system represents a different engineering approach to the rollup problem.

\subsection{Arbitrum}

Arbitrum is a widely adopted Optimistic~Rollup. Its architecture features an interactive fraud proof mechanism implemented through the Cannon virtual machine, combined with a seven-day challenge period that gives validators time to detect and dispute invalid state transitions~\cite{cannon2023}. The system includes sequencer failover capabilities that allow a decentralized validator set to take over transaction ordering if the centralized sequencer goes offline. Arbitrum hosts an ecosystem of decentralized finance protocols, NFT marketplaces, and gaming applications.

The technical challenge in Arbitrum is the implementation of the interactive dispute protocol. The system must correctly order all parties, execute the bisection protocol on-chain, and determine the correct outcome with minimal gas cost. This engineering complexity is reflected in the multi-year development timeline for the Nitro upgrade.

\subsection{Optimism}

Optimism takes a different strategic approach within the Optimistic~Rollup paradigm. Rather than focusing on a single chain, the project developed the OP Stack, a modular software framework that allows third parties to deploy custom rollup instances sharing a common security model and bridging infrastructure~\cite{opstack2023}. This initiative, known as the Superchain vision, has spawned independent chains including Base, Zora, and Mode. Optimism emphasizes EVM equivalence, meaning existing Ethereum smart contracts can be deployed with minimal modification. The \texttt{OptimisticVerifier} contract borrows structural patterns from Optimism Bedrock, particularly the multi-state finality model that distinguishes between pending, soft-final, and hard-final commitments.

The key engineering contribution from Optimism is the Bedrock architecture, which separates the rollup driver from the execution engine and introduces a fault proof system based on Cannon. The modular design allows different fault proof implementations while maintaining the same bridging semantics.

\subsection{zkSync and StarkNet}
\label{sec:zk-production}

zkSync and StarkNet represent two leading production ZK~Rollups, each pursuing a distinct proof system strategy. zkSync employs a PLONK-based proof system and targets full EVM compatibility, allowing developers to write and deploy contracts in Solidity without learning a new language~\cite{zksync-paper}. StarkNet uses the Cairo programming language and its proprietary STARK proof system~\cite{starknet2024,cairo2021}. STARKs offer the advantage of requiring no trusted setup and scale more efficiently for very large computations, though they produce larger proofs than SNARK counterparts. StarkNet offers a Volition mode that allows users to choose on a per-transaction basis whether their data should be stored on-chain for maximum security or off-chain for lower cost.

The engineering challenges in ZK~Rollups differ from Optimistic systems. Proof generation requires significant computational resources, and the verification电路 must be efficient enough to run on-chain within Ethereum's gas limits. Both systems have invested years in hardware acceleration and proof circuit optimization.

Production systems are optimized for normal-case performance. Documentation and public benchmarks focus on throughput under ideal conditions, while behavior under failure scenarios receives comparatively little attention. This optimization bias toward the happy path leaves a gap in understanding of how these architectures perform when things go wrong.

\section{Gap Analysis and Thesis Contribution}

Drawing together the academic literature and industry implementations, we identify five areas where the current body of knowledge lacks empirical evaluation. Each gap corresponds to a specific research question addressed by this thesis.

The first gap concerns finality time measurement. While the theoretical finality properties of both rollup families are well understood, actual finality latencies under varying Layer~1 conditions have not been systematically measured. The framework addresses this gap by recording the precise block number at which each commitment is submitted and the block number at which it reaches finality, enabling exact measurement across diverse scenarios.

The second gap involves recovery cost quantification. When a sequencer crashes or the Layer~1 connection is temporarily lost, the system must recover and resume normal operation. The time and gas costs associated with this recovery process have not been empirically studied. The adversarial test suite includes dedicated crash-and-restart scenarios for both sequencer types, capturing resources consumed during recovery.

The third gap relates to censorship resistance. The ability of honest users to force-include their transactions when a sequencer attempts to censor them is a critical property discussed primarily in theoretical terms. The framework addresses this gap qualitatively in the security analysis and identifies it as a direction for future work.

The fourth gap is the absence of a controlled head-to-head comparison. Existing benchmarks evaluate a single rollup architecture in isolation, making meaningful comparisons difficult. The framework runs both sequencers against the same Layer~1 node simultaneously, ensuring identical block timing, gas pricing, and network conditions for a direct comparison.

The fifth gap concerns reproducibility. Existing performance data for Layer~2 systems comes from mainnet observations or proprietary benchmarks that other researchers cannot replicate. The framework is open-source and deployable with a single Docker Compose command, allowing any researcher to reproduce experiments, verify results, and extend the evaluation to additional rollup architectures.

This thesis addresses the first, second, fourth, and fifth gaps through the framework design, adversarial test suite, and open-source release. The third gap is explored qualitatively for future empirical investigation.

\section{Summary and Transition to Framework}

The literature reveals a clear pattern: both academic research and industry implementations have made significant progress on rollup architectures, but the empirical comparison under adversarial conditions remains largely unexplored. Production systems optimize for the happy path, and academic work focuses on theoretical properties. The gap between theoretical guarantees and practical resilience is the motivation for this thesis.

The next two chapters present the framework design and implementation that enable this empirical comparison. Chapter~\ref{ch:framework} describes the system architecture, technology stack, and operating principles. Chapter~\ref{ch:implementation} provides detailed implementation of the smart contracts, sequencer processes, and metrics infrastructure. Together, these chapters demonstrate that a fair comparison between ZK and Optimistic Rollups is achievable in a controlled academic setting, and establish the experimental methodology for the evaluation that follows.